\section{Task Cost Model}
\label{sec:task-cost-model}

Section~\ref{sec:problem-statement} states that task profiling should produce a
structured estimate rather than a single scalar load value. This section defines
the cost-side part of that estimate. In this paper, task cost has two primary
output families: active resource demand and bounded worker occupancy. Execution
overhead is modeled as a contribution to one or more of those dimensions rather
than as a third independent output family. Operational risk is not part of task
cost; it is modeled later as the consequence of underestimating,
misclassifying, or mishandling an attempt.

The section proceeds from the unit of analysis to the dimensional model. It
first defines cost at task-attempt granularity. It then separates resource demand
from worker occupancy, introduces a per-dimension decomposition into baseline,
contextual amplification, and additive overhead, defines observed values and
residuals, allows distributional estimates, and closes by identifying the
boundaries of the cost model.

\subsection{Cost as an Attempt-Level Quantity}
\label{sec:cost-attempt-level}

The cost model is defined at task-attempt granularity. A task kind determines
the reusable model structure, a task instance supplies intrinsic work features,
and a task attempt is the concrete execution in which demand, occupancy,
observations, outcomes, and residuals are realized.

Let $a \in \mathcal{A}$ be a task attempt. The attempt executes task instance
$i(a) \in \mathcal{I}$, whose task kind is $k(i(a)) \in \mathcal{K}$. For
readability, write
\[
k_a = k(i(a)).
\]
Let $z_a$ be the canonical feature vector available for attempt $a$ after
feature mapping and normalization. We use two role-specific projections of that
vector:
\[
\mathbf{x}_a = P_X(z_a),
\qquad
\mathbf{c}_a = P_C(z_a),
\]
where $\mathbf{x}_a$ contains intrinsic task features and $\mathbf{c}_a$
contains execution-context features. Intrinsic features describe the work
contained in the task instance, such as payload size, record count, object
count, batch size, operation mode, or domain metadata. Execution-context
features describe the environment and state around the attempt, such as worker
state, node pressure, dependency health, network condition, disk pressure, cache
state, retry state, tenant pressure, or domain-specific state.

The distinction between $k_a$, $\mathbf{x}_a$, and $\mathbf{c}_a$ is central. A
task kind may define a reusable cost structure, but it does not determine one
universal cost value. Two attempts of the same task kind may execute the same
logical operation with similar intrinsic features and still differ because they
observe different execution contexts. Conversely, two attempts may execute under
similar context but differ because their intrinsic features represent different
amounts of work. The cost estimate for an attempt is therefore conditional on
task kind, intrinsic features, and execution context.

The model estimates two attempt-indexed vectors:
\[
\widehat{\mathbf{D}}(a)
=
\left\{
\widehat{D}_r(a)
\mid
r \in \mathcal{R}
\right\},
\qquad
\widehat{\mathbf{U}}(a)
=
\left\{
\widehat{U}_o(a)
\mid
o \in \mathcal{O}
\right\}.
\]
The first vector contains active resource-demand estimates. The second contains
worker-occupancy estimates. The attempt-level cost estimate is the ordered pair
\[
\widehat{\mathbf{C}}(a)
=
\left(
\widehat{\mathbf{D}}(a),
\widehat{\mathbf{U}}(a)
\right).
\]
This notation intentionally treats cost as a structured object rather than as a
single scalar.

A scalar may be derived from $\widehat{\mathbf{C}}(a)$ by a downstream policy,
but such a scalar is not the cost model itself. For example, an admission policy,
ranking policy, alerting rule, or profiling-policy selector may apply a
projection
\[
s_\rho(a)
=
\pi_\rho\!\left(
\widehat{\mathbf{D}}(a),
\widehat{\mathbf{U}}(a)
\right)
\]
under some policy $\rho$. The projection $s_\rho(a)$ is policy-specific: it
depends on the decision being made, the available capacity, and the relative
importance assigned to different dimensions. The cost model preserves the
vector-valued estimate because different dimensions have different units,
different composition rules, and different operational consequences.

\subsection{Resource Demand and Worker Occupancy}
\label{sec:cost-demand-occupancy}

The cost object contains two distinct families of attempt-attributed quantities.
Resource demand describes active work required or consumed by the attempt.
Worker occupancy describes bounded execution capacity held by the attempt over
time. The two families may be correlated during execution, but they are not
interchangeable.

For a resource dimension $r \in \mathcal{R}$, $\widehat{D}_r(a)$ denotes the
estimated active resource demand of attempt $a$. Resource-demand dimensions
include quantities such as CPU time, memory peak, allocation volume, disk reads,
disk writes, network input, network output, or accelerator usage. These
dimensions describe the work associated with the attempt itself rather than the
ambient state of the worker, node, dependency, or queue.

For an occupancy dimension $o \in \mathcal{O}$, $\widehat{U}_o(a)$ denotes the
estimated bounded-capacity occupancy of attempt $a$. Occupancy dimensions
include worker-slot time, lease time, connection hold time, lock hold time,
semaphore hold time, or tenant-level concurrency time. These quantities measure
capacity held rather than active resource work. A dependency-bound attempt may
consume little CPU while still holding a worker slot or dependency connection for
a long duration.

The distinction separates demand, occupancy, and pressure. Demand and occupancy
are estimated or observed for the attempt. Pressure is a property of the
execution context. CPU time is a resource-demand dimension; CPU saturation on
the node is a context factor. Disk write bytes are a resource-demand dimension;
disk queue depth is a context factor. Worker-slot time is occupancy attributed
to an attempt; worker-pool pressure is the surrounding state of the pool.

Waiting and elapsed time require the same care. Dependency health, network
condition, connection-pool pressure, and retry state are context factors.
Dependency wait time is an observed waiting signal. Wall-clock duration is an
elapsed-time outcome. Worker-slot time and connection hold time are occupancy
dimensions only when the corresponding bounded capacities remain held during the
elapsed interval. A long dependency wait therefore does not by itself imply equal
worker occupancy: it increases worker-slot time only if the slot remains held.

Two examples illustrate why both families are needed. A compute-dominant
transformation may have high CPU time and allocation volume but release its
worker slot quickly if the computation is short and does not block on external
dependencies. Conversely, an external-service or replication attempt may have
low CPU demand but high worker-slot time, connection hold time, or lease time
when a dependency is slow, a peer is degraded, or retries introduce backoff.
Active resource demand alone would miss the second cost pattern.

\subsection{Baseline, Amplification, and Additive Overhead}
\label{sec:cost-baseline-amplification-overhead}

The model estimates each cost dimension through three mechanisms: intrinsic
baseline, contextual amplification, and additive overhead. The baseline
represents work implied by the task kind and intrinsic features. The
amplification term represents how execution context modifies that baseline. The
overhead term represents additional work or held capacity that is not well
represented as a multiplier.

Let
\[
\mathcal{Q} = \mathcal{R} \cup \mathcal{O}
\]
be the set of cost dimensions. For $q \in \mathcal{Q}$, let
$\widehat{Y}_q(a)$ denote the estimated value for dimension $q$ during attempt
$a$. If $q \in \mathcal{R}$, then $\widehat{Y}_q(a)$ is a component of
$\widehat{\mathbf{D}}(a)$. If $q \in \mathcal{O}$, then $\widehat{Y}_q(a)$ is a
component of $\widehat{\mathbf{U}}(a)$. The dimension-generic estimate is
\[
\widehat{Y}_q(a)
=
\widehat{B}_q(k_a, \mathbf{x}_a)
\cdot
\widehat{A}_q(k_a, \mathbf{x}_a, \mathbf{c}_a)
+
\widehat{O}_q(k_a, \mathbf{x}_a, \mathbf{c}_a),
\qquad
q \in \mathcal{Q}.
\]
Here $\widehat{B}_q$ is the estimated intrinsic baseline, $\widehat{A}_q$ is the
estimated contextual amplification factor, and $\widehat{O}_q$ is the estimated
additive overhead for dimension $q$.

The terms have explicit dimensional constraints. The estimate
$\widehat{Y}_q(a)$, baseline $\widehat{B}_q$, and overhead $\widehat{O}_q$ are
measured in the unit of dimension $q$. The amplification factor
$\widehat{A}_q$ is dimensionless. In the usual cost setting,
$\widehat{B}_q$ and $\widehat{O}_q$ are non-negative, and
$\widehat{A}_q$ is non-negative. Suppression is represented by an amplification
value below one rather than by negative cost.

The baseline term captures cost attributable to the task kind and intrinsic
features. For a compute transformation, record count and algorithm mode may
determine a baseline for CPU time and allocation volume. For a replication task,
payload size or changed-object count may determine a baseline for disk reads,
network output, and commit work. The baseline is not a constant attached to the
task kind alone; it is evaluated using the intrinsic feature values of the task
instance executed by the attempt.

The amplification term captures how execution context modifies the baseline for
dimension $q$. A degraded peer, unstable network, saturated dependency, high disk
pressure, cold cache, or high retry state may amplify some dimensions. A warm
cache, healthy peer, or local data placement may suppress others. The
amplification term depends on both intrinsic features and context because their
interaction matters. Poor network conditions may have little effect on a small
transfer but dominate a large replication attempt.

The overhead term captures cost that should not be treated as proportional to
the baseline. Examples include setup, deserialization, connection establishment,
lease-extension bookkeeping, acknowledgement, commit, cleanup, and accounting or
profiling overhead. Some overheads are nearly constant for a task kind; others
vary with context or intrinsic features. Modeling them additively prevents the
model from forcing every effect into the baseline-amplification product.

This decomposition is an engineering model, not an independence claim. It does
not assert that intrinsic work and execution context are statistically
independent, nor does it require all dimensions to behave linearly. Its role is
to keep three mechanisms separate: work implied by the task, contextual effects
that modify that work, and additional overhead that is not well represented as a
multiplier. The decomposition is also dimension-specific: one context factor may
amplify one dimension, suppress another, and leave a third unchanged.

\subsection{Residuals and Observed Usage}
\label{sec:cost-residuals-usage}

The model distinguishes estimates from observed values. Before or during
execution, the profiler produces $\widehat{Y}_q(a)$ for a cost dimension $q$.
After execution, instrumentation may provide an observed value $Y_q(a)$ for the
same dimension. The residual is
\[
\epsilon_q(a)
=
Y_q(a) - \widehat{Y}_q(a),
\qquad
q \in \mathcal{Q}.
\]
The residual is defined only when an observed value for that dimension is
available. Absence of a residual does not imply accuracy; it may simply mean the
dimension was not observed.

For $q \in \mathcal{R}$, $Y_q(a)$ is an attempt-attributed observed resource
usage value. For $q \in \mathcal{O}$, $Y_q(a)$ is an attempt-attributed observed
occupancy value. A raw measurement is not automatically such a value. Node-level
CPU utilization may be useful context, but it is not observed CPU usage for a
particular attempt. Wall-clock duration may provide evidence, but it becomes an
occupancy value only for the capacity that remains held during the corresponding
interval.

Residuals are not merely accounting errors. They are signals about the relation
between the model and the workload. A large residual may indicate missing
intrinsic features, stale context, hidden dependency state, an incorrect
task-kind classification, model misspecification, measurement error, or a tail
event. A small residual does not prove causal correctness; it only indicates
that the estimate was close to the observed value for that attempt and
dimension.

Residuals connect the cost model to later uncertainty analysis. If residuals are
large, systematic, or concentrated in particular contexts, the model may need
additional features, a revised graph, a different taxonomy, or a higher
profiling level. This section defines residuals only as differences between
observed and estimated cost dimensions; later sections use them to reason about
confidence, drift, and risk.

\subsection{Distributional Cost Estimates}
\label{sec:cost-distributional-estimates}

Attempt-level cost is often variable even within the same task kind and feature
bucket. A single mean estimate can hide behavior that matters operationally,
especially for tasks with retries, dependency waits, lock contention, cache
misses, or degraded peers. The cost model therefore allows each dimension to be
estimated as a distribution or as a set of summary statistics rather than only
as a point value.

For a cost dimension $q \in \mathcal{Q}$ and quantile level $p \in (0,1)$, write
\[
\widehat{Y}_{q,p}(a)
\]
for a quantile estimate of dimension $q$ for attempt $a$. For example,
$\widehat{Y}_{q,0.50}(a)$ may represent a median estimate and
$\widehat{Y}_{q,0.99}(a)$ may represent a high-tail estimate. The paper does not
require one estimator. A system may use empirical quantiles, histograms,
sketches, parametric distributions, conservative envelopes, or model-specific
prediction intervals.

Distributional estimates are useful because different decisions depend on
different parts of the distribution. A mean CPU-time estimate may be sufficient
for coarse accounting of stable compute tasks. A high-percentile worker-slot
estimate may be more relevant for capacity protection when rare attempts can
hold workers for long periods. A high-percentile connection-hold estimate may be
more relevant for dependency protection than average network bytes.

Distributional estimates also prevent a common modeling error: treating
variability as noise around one expected value. In queue-worker systems,
variability may reflect distinct modes. A task may be fast when cache is warm
and slow when cache is cold. A replication attempt may be short when peers are
healthy and long when retries occur. An external-service attempt may be bimodal
because dependency timeouts create a second latency mode. The model should allow
such behavior to remain visible rather than compress it into one average.

This section does not define entropy, confidence, or risk metrics. It only
states that cost estimates may be distributional. Later sections may derive
uncertainty indicators from residuals, tail ratios, drift, missing features,
multimodality, or insufficient sample size.

\subsection{Boundaries of the Cost Model}
\label{sec:cost-model-boundaries}

The cost model defines what is estimated at attempt granularity:
active resource demand and worker occupancy. It does not define the complete
set of influence domains, task behavior classes, graph topology, raw-observation
mapping, uncertainty metrics, risk model, or profiling-policy selection
algorithm.

This boundary is important. Execution context provides input-side conditions
that may influence cost; it is not itself a cost output. Symptoms and outcomes
may provide evidence about cost; they are not automatically cost dimensions.
Risk consumes cost estimates and consequence information; it is not part of the
cost object. Profiling policy may use cost estimates and may introduce profiling
overhead; the decision to enable instrumentation is not made by the cost model.

The next section introduces influence domains, which assign semantic roles to
the inputs, outputs, symptoms, uncertainty indicators, risks, and profiling
observations used by the rest of the paper.
